\section{Classical Smoothness Classes}\label{sec:appendix:smoothness_classes}

This appendix collects the definitions of four classical smoothness classes used in nonparametric estimation theory, stated for functions $f : \mathbb{R}^d \to \mathbb{R}$.
The central distinction is between \emph{isotropic} classes, which couple smoothness across directions, and \emph{anisotropic} classes, which constrain each direction independently.
The main text works with the anisotropic Lipschitz class (\cref{def:anisotropic}); the definitions below provide context for the broader landscape.

The source of coupling in the isotropic Hölder and Sobolev classes is the multi-index: constraining all derivatives of total order $\abs{\alpha} = s$ jointly forces trade-offs between mixed and pure partial derivatives.
The Nikol'skii and Besov classes avoid this by construction --- they are defined via finite differences applied in each coordinate direction separately, making them intrinsically anisotropic.

Throughout, $\alpha = (\alpha_1, \ldots, \alpha_d)$ is a multi-index with $\abs{\alpha} = \alpha_1 + \cdots + \alpha_d$, and $D^\alpha f = \partial^{\abs{\alpha}} f / \partial x_1^{\alpha_1} \cdots \partial x_d^{\alpha_d}$.


\subsection{Hölder / Lipschitz Classes}\label{sec:appendix:holder}

The Hölder class generalizes Lipschitz continuity to higher orders by requiring the highest-order derivatives to satisfy a Hölder condition.

\begin{definition}[Isotropic Hölder Class]\label{def:holder_iso}
	Let $s > 0$ and write $s = \lfloor s \rfloor + \beta$ where $\lfloor s \rfloor \in \mathbb{N}_0$ and $0 < \beta \leq 1$.
	The isotropic Hölder class with smoothness $s$ and constant $M > 0$ is
	\begin{equation*}
		\mathcal{H}^s(M) \defeq \Bigl\{ f \in C^{\lfloor s \rfloor}(\mathbb{R}^d) \;\Big|\; \abs{D^\alpha f(x) - D^\alpha f(y)} \leq M \norm{x - y}_2^\beta \;\;\forall\, x, y,\; \forall\, \abs{\alpha} = \lfloor s \rfloor \Bigr\}.
	\end{equation*}
\end{definition}

The case $s = 1$ ($\lfloor s \rfloor = 0$, $\beta = 1$) gives the Lipschitz class: $\abs{f(x) - f(y)} \leq M \norm{x - y}_2$ for all $x, y$.

\begin{definition}[Anisotropic Hölder Class]\label{def:holder_aniso}
	Let $\mathbf{s} = (s_1, \ldots, s_d)$ with $s_j > 0$, and let $M_j > 0$ for each $j$.
	Write $s_j = \lfloor s_j \rfloor + \beta_j$ with $0 < \beta_j \leq 1$.
	The anisotropic Hölder class is
	\begin{equation*}
		\mathcal{H}^{\mathbf{s}}(M_1, \ldots, M_d) \defeq \Bigl\{ f \;\Big|\; \abs{D_j^{\lfloor s_j \rfloor} f(x) - D_j^{\lfloor s_j \rfloor} f(x + t e_j)} \leq M_j \abs{t}^{\beta_j} \;\;\forall\, x, t,\; j = 1, \ldots, d \Bigr\},
	\end{equation*}
	where $D_j^k$ denotes the $k$-th partial derivative with respect to coordinate $j$ and $e_j$ is the $j$-th standard basis vector.
\end{definition}

Each coordinate direction has its own smoothness order $s_j$ and constant $M_j$.
Variation in direction $j$ is constrained independently of all other directions.
The anisotropic class used in the main text (\cref{def:anisotropic}) is $\mathcal{H}^{(1,1)}(M_x, M_z)$.

\paragraph{The Imbens--Wager smoothness class.}
\textcite{imbens_optimized_2019} work with a second-order isotropic condition: for $X_i \in \mathbb{R}^k$, they assume
\begin{equation*}
	\norm{\nabla^2 \mu_w(x)}_{\mathrm{op}} \leq B, \qquad w \in \{0, 1\},
\end{equation*}
where $\norm{\cdot}_{\mathrm{op}}$ denotes the operator norm (largest absolute eigenvalue of the Hessian).
This bounds the second directional derivative uniformly: $\abs{D_{\vec{v}}^2 \mu_w(x)} \leq B$ for all unit vectors $\vec{v}$ and all $x$.
The operator norm constraint couples second-order variation across all directions --- curvature budget is fungible.

\paragraph{The IW class as the natural agnostic envelope at order 2.}
The operator norm is rotationally invariant: the eigenvalues of $H$ do not change under rotation of the coordinate system.
This makes the IW class the proper analog of the gradient disk at order 1 --- it is the unique rotationally invariant constraint on the Hessian, and therefore the natural agnostic envelope when the researcher is uncertain about the orientation of the smoothness structure.

The isotropic Hölder class $\mathcal{H}^2(M)$ also imposes coupling, but it is \emph{not} rotationally invariant.
It constrains each row of the Hessian to have $\ell_2$ norm at most $M$: in $d = 2$, $f_{xx}^2 + f_{xz}^2 \leq M^2$ and $f_{xz}^2 + f_{zz}^2 \leq M^2$.
This is a coordinate-dependent condition --- it changes under rotation of the axes.
The IW class is strictly contained in the Hölder class with the same constant:
\begin{equation*}
	\{ f : \norm{\nabla^2 f}_{\mathrm{op}} \leq B \} \subsetneq \mathcal{H}^2(B),
\end{equation*}
since $\norm{H}_{\mathrm{op}} \leq B$ implies each row has $\ell_2$ norm at most $B$ (via $\norm{H e_j}_2 \leq \norm{H}_{\mathrm{op}}$), but the converse fails.
The isotropic Hölder class is therefore a looser, coordinate-dependent envelope that pays more than rotation invariance requires; it is included here for completeness.

At order $s = 1$, the distinction vanishes: the gradient norm constraint $\norm{\nabla f}_2 \leq M$ is both the isotropic Hölder class and the unique rotationally invariant constraint.

\paragraph{The cost of agnosticism.}
The primitive object is a rectangle in derivative space: the DGP has bounded partials with some side lengths.
The researcher knows these side lengths but may not know which side aligns with which coordinate axis, and may not know whether or how the partials are coupled.
To cover all possible orientations and coupling structures, the researcher takes the rotationally invariant envelope --- a disk (order 1) or the operator norm ball (order 2) --- with an \emph{inflated} radius.
This is the agnostic class from \cref{sec:overview:epistemology}.

The inflation is the cost of agnosticism:
\begin{itemize}
	\item At order 1: $\mathcal{F}_1^{\mathrm{aniso}}(M, M) \subsetneq \mathcal{F}_1^{\mathrm{iso}}(M\sqrt{2})$.
	The constant inflates from $M$ to $M\sqrt{2}$.
	\item At order 2: the anisotropic primitive $(B_x, B_z)$ inflates to an operator norm bound of $\sqrt{B_x^2 + B_z^2}$.
\end{itemize}
Both the first-order isotropic class and the IW class are the rotationally invariant agnostic envelopes at their respective orders, and both pay the same type of inflation cost.
The inflated constant gives the adversary more room, increasing worst-case bias and minimax risk.

Any information the researcher has about the coupling structure --- or about which direction is smoother --- allows moving from the inflated agnostic class back to the tighter primitive, reducing the constant and improving inference.
This is the central message of \cref{sec:overview:epistemology}: the coupled class is the price of ignorance, and resolving that ignorance strictly helps.

\paragraph{Order-independent structure.}
The table below shows the isotropic and anisotropic classes at orders 1 and 2, with the estimator naturally paired with each order:
\begin{center}
\begin{tabular}{lccc}
	\toprule
	& \textbf{Isotropic (coupled)} & \textbf{Anisotropic (separated)} & \textbf{Estimator} \\
	\midrule
	\textbf{Order 1} & $\norm{\nabla f}_2 \leq M$ & $\abs{\partial f / \partial x} \leq M_x$, $\abs{\partial f / \partial z} \leq M_z$ & Local constant \\
	\textbf{Order 2} & $\norm{\nabla^2 f}_{\mathrm{op}} \leq B$ & $\abs{\partial^2 f / \partial x^2} \leq B_x$, $\abs{\partial^2 f / \partial z^2} \leq B_z$ & Local linear \\
	\bottomrule
\end{tabular}
\end{center}
At order 1, the isotropic constraint forces the gradient into a disk; at order 2, it forces the Hessian eigenvalues into an interval.
In both cases, the curvature budget is fungible across directions, and covariate adjustment cannot reduce worst-case bias.
The anisotropic version fixes each direction's budget independently, enabling covariate adjustment at any order.

Note that the anisotropic class at order 2 places no constraint on the cross-partial $\partial^2 f / \partial x \partial z$.
This is the analog of the anisotropic Lipschitz class placing independent bounds on $\partial f / \partial x$ and $\partial f / \partial z$ without coupling them.


\subsection{Sobolev Classes}\label{sec:appendix:sobolev}

Sobolev classes replace the pointwise (sup-norm) control of Hölder classes with an integrated ($L^p$-norm) control.

\begin{definition}[Isotropic Sobolev Class]\label{def:sobolev_iso}
	Let $s \in \mathbb{N}$ and $1 \leq q \leq \infty$.
	The isotropic Sobolev class is
	\begin{equation*}
		W^{s,q}(M) \defeq \Bigl\{ f \in L^q(\mathbb{R}^d) \;\Big|\; \norm{D^\alpha f}_{L^q} \leq M \;\;\forall\, \abs{\alpha} = s \Bigr\}.
	\end{equation*}
\end{definition}

\begin{definition}[Anisotropic Sobolev Class]\label{def:sobolev_aniso}
	Let $\mathbf{s} = (s_1, \ldots, s_d)$ with $s_j \in \mathbb{N}$, and let $1 \leq q \leq \infty$.
	The anisotropic Sobolev class is
	\begin{equation*}
		W^{\mathbf{s}, q}(M_1, \ldots, M_d) \defeq \Bigl\{ f \in L^q(\mathbb{R}^d) \;\Big|\; \norm{D_j^{s_j} f}_{L^q} \leq M_j,\;\; j = 1, \ldots, d \Bigr\}.
	\end{equation*}
\end{definition}

The isotropic class constrains all partial derivatives of total order $\abs{\alpha} = s$ jointly with one constant --- including mixed partials like $\partial^2 f / \partial x \partial z$ --- coupling smoothness across directions through the multi-index.
The anisotropic class constrains only pure partials $D_j^{s_j} f$ in each direction, with no coupling.
Sobolev classes are natural for $L^2$ estimation (mean integrated squared error); Hölder classes are natural for pointwise estimation and worst-case bias bounds.


\subsection{Nikol'skii Classes}\label{sec:appendix:nikolskii}

Nikol'skii classes \parencite{nikolskii_approximation_1975} are a variant of Sobolev classes that use $L^p$-norms of finite differences rather than derivatives, allowing non-integer smoothness orders.

\begin{definition}[Anisotropic Nikol'skii Class]\label{def:nikolskii}
	Let $\mathbf{s} = (s_1, \ldots, s_d)$ with $s_j > 0$, and let $1 \leq q \leq \infty$.
	The $r$-th difference operator in direction $j$ with step $t$ is
	\begin{equation*}
		\Delta_{j,t}^r f(x) \defeq \sum_{k=0}^{r} \binom{r}{k} (-1)^{r-k} f(x + k t e_j).
	\end{equation*}
	The anisotropic Nikol'skii class is
	\begin{equation*}
		H_q^{\mathbf{s}}(M_1, \ldots, M_d) \defeq \Bigl\{ f \in L^q(\mathbb{R}^d) \;\Big|\; \norm{\Delta_{j,t}^{\lceil s_j \rceil} f}_{L^q} \leq M_j \abs{t}^{s_j} \;\;\forall\, t \in \mathbb{R},\; j = 1, \ldots, d \Bigr\}.
	\end{equation*}
\end{definition}

Nikol'skii classes are intrinsically anisotropic: each coordinate direction has its own smoothness order $s_j$.
When $q = \infty$, the Nikol'skii class coincides with the Hölder class for integer smoothness orders.
The isotropic version is obtained by setting $s_1 = \cdots = s_d$ and $M_1 = \cdots = M_d$.


\subsection{Besov Classes}\label{sec:appendix:besov}

Besov classes refine Nikol'skii classes by adding a secondary integrability parameter that controls how the modulus of smoothness integrates across scales.

\begin{definition}[Anisotropic Besov Class]\label{def:besov}
	Let $\mathbf{s} = (s_1, \ldots, s_d)$ with $s_j > 0$, and let $1 \leq q, r \leq \infty$.
	The anisotropic Besov class $B_{q,r}^{\mathbf{s}}$ consists of functions $f \in L^q(\mathbb{R}^d)$ such that for each $j = 1, \ldots, d$:
	\begin{equation*}
		\norm{f}_{B_{q,r}^{s_j}(j)} \defeq \begin{cases}
			\displaystyle \left( \int_0^\infty \left[ \abs{t}^{-s_j} \norm{\Delta_{j,t}^{\lceil s_j \rceil} f}_{L^q} \right]^r \frac{dt}{\abs{t}} \right)^{1/r} & \text{if } r < \infty, \\[10pt]
			\displaystyle \sup_{t \neq 0}\; \abs{t}^{-s_j} \norm{\Delta_{j,t}^{\lceil s_j \rceil} f}_{L^q} & \text{if } r = \infty,
		\end{cases}
	\end{equation*}
	is finite and bounded by $M_j$.
\end{definition}

The parameter $r$ controls the integration of the modulus of smoothness across scales: $r = \infty$ recovers the Nikol'skii class, while smaller $r$ imposes more uniform smoothness across scales.
For the isotropic version ($s_1 = \cdots = s_d$), mixed-smoothness Besov classes additionally allow cross-directional smoothness interactions.


\subsection{Relationships and Embeddings}\label{sec:appendix:embeddings}

The four classes are related by the following embeddings (stated for the isotropic case with matching parameters):

\begin{center}
\begin{tabular}{lcl}
	$\mathcal{H}^s(M)$ & $\hookrightarrow$ & $H_\infty^s(M)$ \quad (Hölder $\subseteq$ Nikol'skii with $q = \infty$) \\[3pt]
	$B_{\infty, \infty}^s$ & $=$ & $H_\infty^s$ \quad (Nikol'skii $=$ Besov with $r = \infty$, $q = \infty$) \\[3pt]
	$W^{s,q}(M)$ & $\hookrightarrow$ & $B_{q,\infty}^s(M) = H_q^s(M)$ \quad (Sobolev $\subseteq$ Nikol'skii) \\[3pt]
	$B_{q,r_1}^s$ & $\hookrightarrow$ & $B_{q,r_2}^s$ \quad for $r_1 \leq r_2$ \quad (smaller $r$ is more restrictive)
\end{tabular}
\end{center}

In the anisotropic case, each of these embeddings holds coordinate-by-coordinate.

\paragraph{Minimax rates.}
\textcite{stone_optimal_1980} shows that the minimax rate for estimating $f$ at a point under anisotropic smoothness $\mathbf{s} = (s_1, \ldots, s_d)$ is
\begin{equation}\label{eq:stone_rate}
	n^{-2\bar{s}/(2\bar{s}+1)}, \qquad \text{where } \bar{s} = \left(\frac{1}{d} \sum_{j=1}^d \frac{1}{s_j}\right)^{-1} \text{ is the harmonic mean.}
\end{equation}
The minimax optimal estimator uses direction-dependent bandwidths $h_j \asymp n^{-1/(2\bar{s}+1)} \cdot s_j / \bar{s}$, forming a rectangular kernel neighborhood.
Under the isotropic class ($s_1 = \cdots = s_d = s$), this reduces to $n^{-2s/(2s+d)}$ with a spherical neighborhood.

\paragraph{Relevance for the present paper.}
The main text works with anisotropic Lipschitz smoothness ($s_j = 1$ for all $j$, different constants $M_j$) applied to the RDD boundary functional rather than pointwise function estimation.
The choice of Lipschitz (order 1) is driven by the local constant estimator.
Higher-order analogs (using local polynomial estimators of order $p$) would work with anisotropic Hölder classes of order $p + 1$.
The key insight --- that separating smoothness parameters across directions enables covariate adjustment while coupling them prevents it --- applies at all orders.
